\section{Trace transfer and spectral rounding}
\label{ext:sec:rounding}

We approximate trace values by ranks in two steps. First, a positive
function of an element that vanishes near zero is transferred into a
small hereditary support without changing any trace value. Integer
spectral cuts then approximate that value with an error bounded by the
rank of the support. The transferred element may have large norm; the
error will be independent of that norm.

\begin{lemma}[Trace transfer]
\label{ext:lem:transfer}
Let $B$ be a simple stable $C^*$-algebra, let $a\in B_+$ be a
contraction, and let $0\ne h\in B_+$. Suppose that
$r:[0,1]\to[0,\infty)$ is continuous and vanishes on $[0,\delta]$
for some $0<\delta<1$. There is $c\in\her_B(h)_+$ such that
\[
 \tau(c)=\tau(r(a))
\]
for every lower-semicontinuous extended trace $\tau$ on $B$.
\end{lemma}

\begin{proof}
Fullness of $h$ and Lemma~\ref{lem:pre:cut-placement} give
$[(a-\delta/4)_+]\le m[h]$ for some finite $m\ge1$.
Apply that lemma in $M_m(B)$, with cut parameter
$\epsilon=\delta/4$ and comparison error less than $\delta/4$,
to the source $\operatorname{diag}((a-\delta/4)_+,0,\ldots,0)$
and target $H=\operatorname{diag}(h,\ldots,h)$.
The resulting placement matrix has only its first column nonzero,
because all other diagonal entries of its initial squared modulus
are zero. Thus there is a column $v=(v_i)_{i=1}^m$ such that
\[
 v^*v=(a-\delta/2)_+,\qquad
 vv^*\in M_m(\her_B(h)).
\]
Define
\[
 g(t)=\left(\frac{r(t)}{(t-\delta/2)_+}\right)^{1/2},
\]
setting it equal to zero where the denominator vanishes.
It is continuous and vanishes near zero because $r$ vanishes on
$[0,\delta]$. For $z_i=v_i g(a)$, we have
\[
 \sum_{i=1}^m z_i^*z_i=r(a),\qquad
 c:=\sum_{i=1}^m z_i z_i^*\in\her_B(h).
\]
For the containment, use
$z_i z_i^*\le\|g\|^2v_i v_i^*$ and heredity.
Additivity and traciality give the asserted identity for every
extended trace, with the same $c$.
\end{proof}

\begin{lemma}[Integer spectral rounding]
\label{ext:lem:rounding}
Let $B$ be a stable $C^*$-algebra, $c\in B_+$, and let $L$ be an
integer greater than $\|c\|$. Put
\[
 w=\sum_{i=1}^L[(c-i)_+]\in\Cu(B).
\]
For every lower-semicontinuous extended trace $\tau$, writing
$\varphi=d_\tau$, one has
\begin{equation}
\label{ext:eq:rounding}
 \varphi(w)\le\tau(c)\le\varphi(w)+d_\tau(c).
\end{equation}
\end{lemma}

\begin{proof}
For $q(t)=\min(t,1)$ on $[0,\infty)$, scalar functional calculus
gives the exact identity
\begin{equation}
\label{ext:eq:staircase}
 c=\sum_{i=1}^L q((c-(i-1))_+).
\end{equation}
Set $R=\max(1,\|c\|)$. For every integer $m\ge1$,
\[
 \big((c-i)_+/R\big)^{1/m}\le q((c-(i-1))_+).
\]
Indeed, below the spectral threshold $i$ the left side is zero,
and above it the right side is one. Consequently,
\[
 d_\tau((c-i)_+)
 \le\tau(q((c-(i-1))_+))
 \le d_\tau((c-(i-1))_+).
\]
The last inequality uses Equation~\eqref{eq:pre:trace-rank} and
hereditary containment. Summing and using $(c-L)_+=0$ proves
Equation~\eqref{ext:eq:rounding}. Only the first rank $d_\tau(c)$
remains as error; the number $L$ of spectral levels does not occur
in the error bound.
\end{proof}

We now combine the two lemmas. If $M[h]\le[e]$ and $c$ is obtained
by trace transfer into $\her(h)$, then $M[c]\le[e]$.
For a functional with $\varphi([e])<\infty$, rounding therefore
approximates $\tau(r(a))$ to within $\varphi([e])/M$.
A single cut-supported function does not recover the full rank of an
arbitrary element. In the next section we write that rank as a sum of
such trace values and choose the support ranks so that the errors are
summable. Neither the norm of $c$ nor the number of its spectral levels
has to be bounded uniformly.
